
\begin{comment}
    

\subsection {Related Works of GraphRAG}


\textbf{Retrieval-Augmented Generation (RAG).} 
Retrieval-Augmented Generation (RAG) addresses the limitations of LLM-generated content by incorporating knowledge from external databases, thereby improving factual accuracy and contextual relevance in the outputs.
Early RAG frameworks relies on basic retrieval mechanisms, such as chunking documents and applying text embeddings with cosine similarity for top-$k$ retrieval~\cite{lewis2021retrieval}.
These naive approaches, however, often retrieved noisy or irrelevant information, leading to suboptimal reasoning and degraded performance~\cite{HyDE}.

To address these shortcomings, advanced RAG frameworks incorporate more refined retrieval mechanisms to improve document filtering and ranking. 
 Focusing on improving factual accuracy and contextual relevance, these methods utilize pre-retrieval~\cite{StepBackRAG,Rewriting,HyDE} and post-retrieval~\cite{G-RAG,DPA-RAG,R4} mechanisms to enhance the quality of retrieved information. {\color{red} \bf Be more specific on techniques rather than keep repeating improving performance, relevance and accuracy.}

Recent advancements have shifted toward modular RAG approaches~\cite{surveyTJFD}, such as LlamaIndex~\cite{LlamaIndex} and LangChain~\cite{LangChain}, which decompose the RAG process into distinct modules (e.g. rewrite, retrieve, rerank, fusion, memory???) for more flexible and efficient retrieval and generation. 
These frameworks not only improve reasoning capabilities and factual accuracy, but also underscore the potential of modular design for language understanding and generation tasks. 
For instance, by separately fine-tuning retrieval and generation stages, approaches such as those by \citet{yu2023generateretrievelargelanguage} allow for better control over the RAG process. 
Despite these advancements, RAG still encounters challenges with noise and irrelevant information retrieval~\citep{GraphRAG}.

\textbf{Existing GraphRAG Implementations (or Instances).}
In GraphRAG, graphs serve as substitutes for document sets in RAG retrieval, allowing the extraction of accurate ``reasoning paths'' tailored to the query context. This approach enhances LLM reasoning by leveraging structured knowledge.  
Graphs provide structured representations of knowledge, that aid in improving the relevance and accuracy of retrieved reasoning paths.
In addition, graph-based searching and ranking algorithms effectively reduce noise, addressing the redundancy issues commonly associated with document-based retrieval methods.
While the study of GraphRAG is in its early stages, with only limited implementations or instances available, as summarized in Table~\ref{tab:Current_IMP}, its potential for broader applications across diverse domains is promising, paving the road for more effective RAG systems.

For the core retrieval phase in GraphRAG, current instances often incorporate techniques from the research line of knowledge-base question-answering (KBQA), utilizing both information retrieval (IR)~\cite{KV-Mem,EmbedKGQA,NSM,TransferNet,GraftNet,PullNet} and semantic parsing (SP)~\cite{QGG,HGNet,Program-Transfer,UniKGQA,DecAF} methods to identify relevant subgraphs or reasoning paths for question answering.
For example, GCR~\cite{GCR} utilizes LLMs and beam search to construct precise reasoning paths that are closely aligned with user queries, improving the relevance and quality of retrieved reasoning paths. Similarly, GSR~\cite{LessIsMore} incorporates small-scale specialized models, such as T5, as alternatives to LLMs to optimize the retrieval performance. 
RoG~\cite{RoG}, GSR\cite{LessIsMore} and GCR~\cite{GCR} apply Personalized PageRank (PPR) for retrieving relevant query-specific subgraphs, whereas KELP~\cite{ref:kelp} uses a fine-tuned BERT model~\cite{bert} to refine reasoning paths for better generation quality. 

Additionally, many instances incorporate LLMs or fine-tuned LLMs into the retrieval phase. 
For example, RoG~\cite{RoG} leverages fine-tuned LLMs to identify key entities and relationships, while ToG~\cite{ToG} and DoG~\cite{DoG} leverage vanilla LLMs to retrieve relevant reasoning paths.

In summary, these methods showcase the variety of retrieval techniques in GraphRAG, from graph-based algorithms to neural network models. By integrating structured graph knowledge with advanced retrieval, they enhance accuracy and contextual relevance in question answering, paving the way for overcoming the limitations of document-based systems.

\end{comment}